%-------------------------
% Rover Resume - Base Template
% Link: https://github.com/subidit/rover-resume
%
% Shows code for various formatting options for different resume sections.
% Education and Projects have single-line headers; while Experience uses double-line.
% Some formatting codes are kept inline; consider \newcommand{cmd}{def}.
% Excludes hyperref and icons for readability; MVP version.
% Explore other templates for more options.
% Mix and match as desired. Be consistent with headers and sub-headers.
%------------------------

\documentclass[11pt]{article} % fontsize 10pt/11pt/12pt

\usepackage[margin=1in, a4paper]{geometry}
\setcounter{secnumdepth}{0} % remove section numbering
\usepackage{titlesec}
\titlespacing{\subsection}{0pt}{*0}{*0} % remove vertical spacing above and below
\titlespacing{\subsubsection}{0pt}{*0}{*0}
\titleformat{\section}{\large\bfseries\uppercase}{}{}{}[\titlerule]
\titleformat*{\subsubsection}{\large\itshape}
\usepackage{enumitem}
\setlist[itemize]{noitemsep,left=0pt .. \parindent}
\pagestyle{empty} % remove page number

% Bibliographies live in cv/assets/, split by kind so each section draws from
% its own file. doi=true keeps the DOI visible: APS's new DOI format cannot be
% guessed from volume/page, so it is the only manual path back to the paper.
\usepackage[backend=biber, style=trad-plain, doi=true, url=false, isbn=false,
            giveninits=false, maxbibnames=99]{biblatex}
\addbibresource{../assets/publications.bib}
\addbibresource{../assets/oral.bib}
\addbibresource{../assets/poster.bib}


\begin{document}

\begin{center}
	\begin{minipage}{0.5\textwidth}
		{\Huge\bfseries
			Sota Shimozono % Name Here
		} \\ \medskip
	\end{minipage} \hfill
	\begin{minipage}{0.4\textwidth}
		\raggedleft
		Email: shimozono-sota631@g.ecc.u-tokyo.ac.jp \\
		LinkedIn: linkedin.com/in/sota-shimozono-07a696263 \\
		GitHub: github.com/sotashimozono
	\end{minipage}
\end{center}

\section{Education}
%=================%
\subsection{University of Tokyo $|$ {\normalfont\itshape Master of Science} \hfill Apr 2025 -- Mar 2027 (Expected)}
\begin{itemize}
	\item Department of Basic Science, Graduate School of Arts and Sciences
	\item Research topic: One-dimensional quantum many-body systems
	\item Related Coursework: Computational Physics
	\item Intending to pursue a Ph.D. program
\end{itemize}
\subsection{University of Tokyo $|$ {\normalfont\itshape Bachelor of Arts and Sciences} \hfill Apr 2021 -- Mar 2025}
\begin{itemize}
	\item Department of Integrated Sciences, Subcourse in Mathematical Natural Sciences
	\item Related Coursework: Mathematics, Physics, Computational Physics
\end{itemize}

\section{Fellowships}
%====================%
\begin{enumerate}[label=\null, left=0pt..0pt, itemsep=2pt]
	\item {\bfseries WINGS-ABC Fellow}, The University of Tokyo \hfill \mbox{2025 -- 2030} \newline
	World-leading INnovative Graduate Study Program of Advanced Basic Science Course; selected by written application and interview, with a stipend.
\end{enumerate}

\section{Publications}
%=====================%
% series/resume keeps one running number across Publications and
% Presentations, so an entry can be cited unambiguously as [n].
\begin{enumerate}[series=cvrefs, label={[\arabic*]}, left=0pt..\parindent, itemsep=2pt]
	\item \fullcite{shimozono2026envmpo} \newline
	Recovers thermodynamic-limit bulk behavior by attaching environment operators to both ends of a finite-size calculation, without an infinite-size algorithm.
\end{enumerate}

\section{Presentations}
%======================%
\subsection{Orals}
\begin{enumerate}[resume=cvrefs, label={[\arabic*]}, left=0pt..\parindent, itemsep=2pt]
	\item \fullcite{shimozono2025jps_hiroshima}
	\item \fullcite{shimozono2025jps_online}
\end{enumerate}

\subsection{Posters}
\begin{enumerate}[resume=cvrefs, label={[\arabic*]}, left=0pt..\parindent, itemsep=2pt]
	\item \fullcite{shimozono2025tnsaa}
	\item \fullcite{shimozono2025ttqm}
\end{enumerate}

\section{Experience}
%=================%
\subsection{The University of Tokyo, College of Arts and Sciences \hfill Komaba, Tokyo}
\subsubsection{Teaching Assistant \hfill Apr 2024 -- Present}
\begin{itemize}
	\item Twelve assignments across six subjects and six semesters, peaking at four concurrent courses in 2025A with over 40 contact sessions in that term alone:
	\begin{itemize}
		\item Introductory Physics Experiments --- 2024S, 2024A, 2025S, 2025A
		\item Electromagnetism --- 2024A (two sections), 2026A (scheduled)
		\item Oscillations and Waves --- 2025S, 2025A
		\item Quantum Mechanics --- 2025A
		\item Exercises in Mathematical Sciences --- 2025A
		\item Introduction to Physics --- 2026S
	\end{itemize}
	\item Developed \texttt{BlochSpins.jl}, an open-source Julia package for simulating and visualizing quantum states on the Bloch sphere, to address common student inquiries in Quantum Mechanics.
\end{itemize}

\subsection{NPO Mochikoma $|$ KOMAD student community space \hfill Komaba, Tokyo}
\subsubsection{Director \hfill 2024 -- Present} % TODO: confirm the year the directorship started
\subsubsection{Steering Committee Member \hfill Mar 2023 -- Nov 2024}
\begin{itemize}
	\item Management of the space and community.
	\item Organized the community barbecue in 2023 and 2024.
	\item Changed the operating policy of the space in Nov 2024, re-establishing membership on new terms.
\end{itemize}

\subsection{KUREHA, University of Tokyo Tea Society \hfill Komaba, Tokyo}
\subsubsection{President \hfill Apr 2023 -- Mar 2025}
\subsubsection{Vice President \hfill Apr 2022 -- Mar 2023}
\begin{itemize}
	\item Developed an original milk-tea blend and took it to commercial release.
	\item Ran a pop-up tea cafe at Gallery Cafe NICCO at Shin-Nakano, an off-campus venue (2023).
	\item Joined a tea-leaf sourcing trip to Sri Lanka with WonderLinks (Feb 2025).
\end{itemize}

\subsection{University of Tokyo Frontier Runners \hfill Samani, Hokkaido}
\subsubsection{Volunteer, educational outreach \hfill Aug 2022 -- Aug 2024}
\begin{itemize}
	\item Took part in five field sessions over three years, from summer 2022 to summer 2024.
	\item Supported study and university-entrance preparation for local high-school students.
\end{itemize}

\section{Skills \& Interests}
%===========================%
\begin{description}[itemsep=0pt]
	\item[Technical] Julia / Rust / Python / TypeScript / Lean 4 / \LaTeX{} / Git / GitHub Actions
	\item[Research] Tensor networks (MPS, DMRG, TDVP) / quantum many-body physics / numerical linear algebra / HPC on shared clusters
	\item[Language] Japanese (native) / English (research reading and writing, everyday conversation)
	\item[Interests] Tea / Coffee / Cocktails / Community Building / Open Source Software
\end{description}

\section{Open Source}
%====================%
\begin{itemize}
	\item {\bfseries obsidian-remote-ssh} --- Obsidian plugin for editing remote vaults over SSH/SFTP. Built in TypeScript and Go. Listed in the Obsidian community plugin directory.
	\item {\bfseries doiget} --- CLI and stdio MCP server that resolves DOIs and arXiv IDs to local PDFs through open-access APIs. Built in Rust and served on crates.io and the MCP Registry.
	\item {\bfseries QAtlas.jl} --- Julia package collecting rigorous reference results for quantum many-body systems.
	\item Several smaller Julia packages in the General registry, covering lattice construction, quasicrystals, parameter handling, and CI tooling.
\end{itemize}

\end{document}
